\subsection{Findings Overview}
\label{appendix-findings-table}
\begin{longtable}{|p{3cm}|p{3.5cm}|p{7.5cm}|}
\hline
\rowcolor[HTML]{B3E5FC} 
% Deeper Connection Section
\multicolumn{3}{|c|}{\cellcolor[HTML]{E1F5FE}\textbf{I - Deeper Connection}} \\ \hline
\textbf{Aspect} & \textbf{Keywords} & \textbf{Concepts} \\
\hline
Restorative Effects & Brain fatigue, recovery, attention restoration theory (ART) & Nature’s ability to relieve cognitive overload and stress. Restorative and therapeutic impacts of nature interaction.  \\ \hline

Existential Experiences with Nature & Ego reflection, awe, discomfort & Invoke more existential notions. Nature invokes a sense of awe and scale that shifts human focus from control to reflection and humility.Encounters with nature can provoke discomfort or fear, but these moments of tension are opportunities for deeper reflection, growth, and meaning-making.
\\ \hline

Genius Loci and Place Identity & Spirit of place, unique identity, spatial design, manifestation of place & Places have distinct historical and ecological qualities. Spatial design can amplify and make perceivable the unique essence or identity of a location. \\ \hline

Multi-Sensory Experiences & Smell, vision, sound, touch, taste, movement & Biophilic design should engage multiple senses to deepen human-nature connections. It’s not just about visual greenery but creating environments that feel alive. \\ \hline


Participation and Engagement & Participatory planning, maintenance, ownership, community & Engaging users in the care and maintenance of nature fosters deeper personal and collective connections with the natural elements in their environment. \\ \hline

Awareness and Respect & agency, cultural gap & Recognising nature as an active agent rather than a passive setting. Also, awareness of the benefits and importance of nature leads to respect and stewardship. Bridging cultural gaps in nature perception is key. \\ \hline

Inter-Subjectivity & Non-human beings, habitats, Active Ecosystems, inter-subjectivity relationships &  Developing reciprocal relationships between humans and other living beings. Nature in design should reflect living, thriving ecosystems, not artificial or superficial elements.  \\ \hline

% Existing Design Limitation Section
\multicolumn{3}{|c|}{\cellcolor[HTML]{C8E6C9}\textbf{G - Existing Design Limitation}} \\ \hline
\textbf{Aspect} & \textbf{Keywords} & \textbf{Concepts} \\
\hline
Limited multi-sensory engagement & Visual perception, image-driven, first connection, visual impact, indirect forms & Visual perception dominates biophilic design strategies but lacks depth in engaging other sensory dimensions. Current designs fail to incorporate diverse sensory stimuli. Multi-sensory layering could add depth and intensify connections.\\ \hline


Scale and Connectivity & Vast, smaller scale, ecological system, leave room for nature, connect to the ground & Overly vast spaces dilute engagement. Maintaining ground connectivity is essential for impactful biophilic integration. \\ \hline

Superficial Approaches to Nature & Fake plastic plants, indirect forms, cheapest way, no interactions, Friction, reality, absurdity, minimal performance, weak connections, Perception and Engagement  & Superficial elements prioritize aesthetics over authentic, interactive, or ecological value. Designs fail to create friction or challenge perceptions, offering only shallow connections to nature.\\ \hline


% Sense Section
\multicolumn{3}{|c|}{\cellcolor[HTML]{F1F8E9}\textbf{G - Sense}} \\ \hline
\textbf{Aspect} & \textbf{Keywords} & \textbf{Concepts} \\
\hline
Smell, sound, taste, and touch & Smell, sound, touch, humidity, sensory experience, multi-sensory & Smell is complex and often suppressed in current designs due to concerns about unpleasant odors. Sound remains underutilized as current practices focus on localization and minimization. Tactile experiences complement visual ones but are rarely prioritized. Taste is challenging and almost never considered in built environments. Collectively and individually, these senses can contribute to constructing the meaning of a space. \\ \hline

Movement & Flow, motion, kinetic design, spatial interaction & Movement as a 6th sense adds a dynamic element to biophilic design, engaging users physically and emotionally. Designing in a way that allows for active movement  of users  that encourage exploration can create a sense of vitality and connection to natural rhythms. \\ \hline

Sensory Layering and Gradual Engagement & Layers, gradual engagement, sensory mapping & Sensory layering can be employed to guide users through various sensory dimensions, gradually enhancing their engagement. Workshops and sensory mapping exercises emphasize using more than one sense at a time, demonstrating the value of multi-dimensional sensory design. \\ \hline

Designing for Interaction and Response & Reaction, interaction, response, mediation of presence & Spaces designed to react to human interaction, such as lighting changes when grass is touched or walls mediating presence, offer dynamic and engaging experiences. Interactive biophilic design blurs the line between passive observation and active participation, creating a sense of agency and connection. \\ \hline

Complexity and Intelligence in Nature & Ecosystem, biomes, communication, control, rejection & Nature’s complexity, including microbial processes and ecosystem dynamics, enriches biophilic designs by evoking curiosity and fostering deeper human-environment connections. Engaging with discomfort, unpredictability, or innate fear in nature adds depth and meaning to these experiences. However, the need for control over sensory elements, particularly in indoor environments, poses challenges to incorporating the untamed aspects of nature. \\ \hline

% Challenges and Feasibility Section
\multicolumn{3}{|c|}{\cellcolor[HTML]{E3F2FD}\textbf{G - Challenges and Feasibility}} \\ \hline
\textbf{Aspect} & \textbf{Keywords} & \textbf{Concepts} \\
\hline
Regulatory, Economic, and Technical Constraints & Tropical rainforest species, nurseries, tender processes, technical limitations & Availability of specific plant species and reliance on nurseries limit biophilic design choices. Tender processes create barriers for unconventional or innovative ideas, requiring alignment with existing capabilities and resources. Biophilic design must demonstrate tangible economic benefits to gain acceptance in the market. \\ \hline


Morphological Constraints & Historical city centers, fixed morphology,  & Strict preservation laws in historical cityscapes complicate the integration of nature-inspired elements. Achieving a balance between maintaining cultural heritage and improving urban microclimates is a persistent challenge. \\
\hline

Liability and Maintenance Concerns & Risk aversion, liability, water leakage, maintenance costs & Concerns over long-term maintenance and liability deter stakeholders from adopting biophilic solutions. High upfront costs and perceived risks create resistance among project developers. \\
\hline

Knowledge and Data Limitations & Sound data, sensory mapping, location-specific sound management, noise cancellation & The limited ability to utilize sensory data (e.g., sound) creatively restricts biophilic innovation. Current practices often rely on simpler solutions like noise cancellation rather than exploring sensory complexity. \\
\hline

Design Trade-Offs & Vertical buildings, ground-level connections, horizontal design, structural limitations & Urban expansion through vertical architecture often sacrifices ground-level, nature-integrated spaces. Effective biophilic design requires reconciling the trade-offs between vertical and horizontal approaches. \\
\hline

% Criticisms Section
\multicolumn{3}{|c|}{\cellcolor[HTML]{FFF3E0}\textbf{G - Criticisms}} \\ \hline
\textbf{Aspect} & \textbf{Keywords} & \textbf{Concepts} \\
\hline
Tokenistic Approaches & Pure decoration, greenwashing, fake plants, artificial green walls, superficial & Biophilic design is often reduced to superficial aesthetics, such as fake plants or green walls with no real ecological or sensory benefit. Greenwashing practices undermine the authenticity and impact of biophilic principles. \\
\hline

Misalignment with Biophilic Philosophy & Energy-consuming, waste-consuming, against biophilic philosophy, respect for nature & Over-engineered or resource-intensive solutions contradict the principles of ecological respect and sustainability inherent in biophilic design. Designs that prioritize form over function fail to align with biophilic goals of fostering deeper ecological awareness. Treating plants and greenery as mere decorative elements rather than active participants in the ecosystem diminishes their role in biophilic design. This reductionist view limits opportunities for meaningful human-nature interactions.\\
\hline


Idealized and Unrealistic Visions of Nature & Idealized visions, overgrown cities, exaggerated greenery, utopia & Representations of future cities with extreme greenery are often impractical and fail to consider the complexities of integrating nature into urban spaces. These visions create unrealistic expectations and detract from achievable biophilic goals. \\
\hline

Lack of Ecological Coherence & Natural tools, planetary resources, coherence, real nature & Biophilic designs often lack ecological coherence, using synthetic or isolated elements that fail to replicate real natural processes. This approach undermines the goal of fostering authentic connections with nature. \\
\hline

% Emerging Materials and Tech Section
\multicolumn{3}{|c|}{\cellcolor[HTML]{E3F2FD}\textbf{G - Role of Technolgy}} \\ \hline
\textbf{Aspect} & \textbf{Keywords} & \textbf{Concepts} \\
\hline

Climate-Responsive Technologies & Semi-passive coatings, smart glass, adaptive blinds, seasonal transformations & Climate-responsive technologies like smart glass and semi-passive coatings dynamically adapt to environmental conditions, balancing comfort and energy efficiency. Seasonal adaptations in materials further enhance sustainability and occupant well-being. \\ \hline

Reactive Biophilic Booths & Touch-responsive surfaces, reactive spaces, light, sound, smell, tactile feedback, green capsules & Reactive biophilic booths integrate interactive elements like touch-responsive walls and reactive gardens to engage users with their surroundings. These immersive designs simulate natural interactions, fostering connection and agency. Multisensory integration, including light, sound, smell, and tactile feedback, enhances user experience, while innovations like green capsules provide restorative, multi-dimensional environments. \\ \hline

Health-Centred Multi-Sensory Design & Biodomes, exploration, smell, circadian rhythm, sensory optimization & Biophilic designs that incorporate smell, sound, and touch are instrumental in creating multi-sensory interactions that promote physical and psychological well-being. Innovations such as circadian rhythm-adjusted lighting systems help align indoor environments with natural biological cycles, enhancing health outcomes. Biodomes, as immersive environments, integrate diverse sensory elements like textures, smells, and exploratory pathways to foster deeper connections with nature while supporting overall health and wellness. \\
\hline

Behavioral Prompts and Nudging & Reminders, prompts, engagement cues, nature interactions & Technology supports biophilic design by using reminders and prompts to nudge occupants toward engaging with natural elements. Subtle cues like digital notifications or environmental changes encourage interaction and enhance well-being. \\ \hline

Educational Technologies & QR codes, sensors, green routes, educational layers, user engagement & QR codes and sensors enhance engagement by providing real-time plant information. Interactive routes and post-design interventions integrate learning into biophilic spaces, making them more engaging and alive. \\ 
\hline

Hybrid Craft-Tech Synergy & Craft-tech synergy, dynamic behaviors, 3D printing, computational approaches & Combining traditional craft methods with advanced computational techniques, such as 3D printing and algorithmic design, results in biophilic environments that simulate natural patterns and behaviors. \\ \hline

Scalable Biophilic Urbanism & City-scale biophilia, urban application, performance-based measurements, ecological processes & Applying biophilic design principles to urban scales involves adapting smaller building-level strategies (e.g., green facades, adaptive coatings) to city-wide ecosystems. Performance-based metrics and regulations can ensure urban biophilic spaces perform effectively, enhancing ecological and human health outcomes.  Material innovations, such as the “Soft Wall,” utilize patterns, textures, and dynamic properties to mediate emotional and spatial connections between users and environments. These materials go beyond aesthetics to serve functional roles, such as enhancing privacy, creating sensory experiences, or fostering communication between spaces.\\ \hline

Integration of AI and Smart Systems & AI interaction, Alexa, mediating presence, direct interaction with environments, subtle sensors, digital tools, TPD-T model & Smart systems like AI and voice-controlled technologies can enable direct interaction with buildings, vegetation, and surroundings, enhancing user engagement. Examples include voice-guided interactions with nature-inspired installations, fostering connections between users and their environments. I and smart systems enhance biophilic experiences through responsive and subtle adjustments, such as light gradients. Digital tools like the TPD-T model enable biophilia-informed urban planning and policy-making.  \\
\hline

Realism in Artificiality and Seamless Integration & Mixed Reality (MR), naturalistic simulation, immersive design, blended environments & Mixed Reality (MR) technologies blur the boundaries between artificial and natural environments by overlaying digital natural elements onto physical spaces. These tools enable immersive biophilic experiences, simulating realistic patterns, textures, and transitions that evoke emotional connections. By integrating physical and digital layers seamlessly, MR fosters a sense of presence and connection to nature, avoiding artificial hard lines or borders that disrupt the immersive experience. \\ \hline


% Values Section
\multicolumn{3}{|c|}{\cellcolor[HTML]{E8F5E9}\textbf{E - Values}} \\ \hline
\textbf{Aspect} & \textbf{Keywords} & \textbf{Concepts} \\
\hline
Inclusivity and Justice & Inclusivity, diversity, accessibility, social justice, spaces of injustice & Biophilic design must create spaces accessible to all social groups, fostering equity and diversity in public and private environments. Design can address societal injustices, especially in marginalized or transient communities, by improving access to natural spaces and fostering a sense of permanence and belonging. Affluence often correlates with access to green spaces, highlighting the need for nature-inclusive design in less affluent areas. \\
\hline
Artistic Expression & Artistic intervention, spatializing practice, storytelling, art and nature & Artistic installations and storytelling can transform biophilic spaces into expressions of community identity and collective memory. Public spaces that reflect diverse perspectives foster a shared sense of place and cultural relevance. The integration of art with natural elements enhances emotional engagement, creating spaces that tell the stories of the people and the land. \\
\hline
Socialisation & Balance, social interaction, control, sensitivity, burnout recovery & Biophilic design should prioritize user control over the environment to balance work-life dynamics and reduce stress. Environments that support personal well-being and encourage social interaction foster emotional resilience. Sensitivity and perception, including gendered differences in sensory engagement, play critical roles in how individuals connect with nature. \\
\hline
Feminism & Sensitivity, human sensor, gendered perception, inner connection & The human body acts as a natural sensor, fostering connections with the environment through sensory perception. Gendered differences in sensitivity, with women being more attuned to biophilic experiences, highlight the importance of designing for diverse needs. Sensory design should amplify this innate human ability to reconnect with nature, promoting ecological awareness and respect. \\
\hline
Community and Collective Responsibility & Social environments, community engagement, participatory design & Participatory planning processes foster ownership and stewardship, encouraging communities to maintain and care for green spaces. Shared spaces for interaction and collaboration can strengthen social bonds and promote collective responsibility for ecological well-being. \\
\hline
Global urgencies  & Global urgencies, scalability, ecological impact & Design practices should address global challenges, from climate change to urbanization, ensuring scalable solutions that balance local needs with broader ecological goals. Even small-scale interventions in biophilic design can create ripple effects, fostering systemic change and promoting sustainability. \\
\hline
% Future Section
\multicolumn{3}{|c|}{\cellcolor[HTML]{FFF8E1}\textbf{V - Future}} \\ \hline
\textbf{Aspect} & \textbf{Keywords} & \textbf{Concepts} \\
\hline

Biodomes and Ecosystem Participation & Ecosystem, immersive, participation, balance & Indoor biodomes can immerse occupants in ecosystems, reminding them of their role within larger natural systems. Spaces should emphasize coexistence and participation rather than domination of nature. ndoor biophilic spaces can serve as mini-biomes or ecological corridors for urban wildlife, enabling dynamic relationships between species. Technology can monitor and manage these systems, balancing populations and providing opportunities for building occupants to observe and engage with wildlife \\ \hline

Multi-Emotional Spaces & Emotional states, body rhythms, tailored environments, gendered needs & Spaces can cater to diverse emotional needs and cycles, such as areas for cocooning or energizing. Technologies that monitor body rhythms and emotions can adapt indoor environments to align with individual well-being and productivity. \\ \hline

Community-Centered Biophilic Design & Community spaces, youth programs, escape, safe spaces & Biophilic design for public buildings, such as libraries or community centers, can offer programming that supports community needs, particularly for youth. These spaces can provide safe, accessible refuges for social interaction and personal escape. \\ \hline

Speculative Smart Nature & Smart nature, invisible walls, immersive biophilic technology & Future designs can integrate invisible walls or speculative technologies to simulate the sensation of being outdoors while remaining sheltered. Sensors and screens could provide restorative and immersive nature-inspired moments indoors. \\ \hline

Infinite Space and Horizons & Transparency, reflections, infinity, horizon views & Transparent materials and reflective surfaces can create a sense of infinite space and horizon indoors, enhancing the perception of being connected to nature on a grand scale. \\ \hline

Dynamic Time and Rhythms & Tides, seasonal change, shadows, flux & Indoor spaces can replicate nature’s dynamic interactions with time, such as changing tides, shifting shadows, and evolving seasons. These elements tie into humans' innate perception of being part of larger natural rhythms. \\ \hline

Zero Waste and Circular Systems & Waste-free, circular processes, natural cycles & Indoor systems can emulate nature’s waste-free, cyclical processes to promote sustainability. These designs can integrate resource efficiency and interconnected ecological frameworks. \\ \hline
\end{longtable}